\section{Introduction}

Dense retrieval is attractive when lexical matching is too brittle, but its deployment cost is not symmetric. In many production settings, document embeddings can be computed offline and indexed once, whereas query encoding lies directly on the online critical path. This difference makes query-side efficiency more operationally important than overall encoder symmetry. The relevant question is therefore not only how to improve a dual encoder in the style of DPR and related bi-encoder systems \cite{karpukhin-etal-2020-dense,reimers-gurevych-2019-sentence}, but how much retrieval quality survives when the document tower remains strong and only the query tower is compressed.

This paper studies that question in a deliberately narrow form. We fix a strong zero-shot teacher, build asymmetric students by truncating the query tower alone, and then test a small sequence of low-cost interventions: pooling, a projection head motivated by contrastive representation learning \cite{chen-etal-2020-simclr}, lightweight distillation \cite{huang-chen-2024-pairdistill}, and a later approximate-nearest-neighbor (ANN) systems comparison implemented with standard FAISS indices \cite{johnson-etal-2021-faiss,malkov-yashunin-2018-hnsw}. The goal is not to propose a new state-of-the-art retriever. The goal is to isolate failure modes and first-order tradeoffs in online-efficient dense retrieval.

The resulting empirical picture is sharper than the original hypothesis suggested. A zero-shot \bgebase\ teacher remains extremely strong on the prepared MS MARCO-derived benchmark. Naive 4-layer and 2-layer query truncation causes catastrophic retrieval collapse. Query-side mean pooling becomes the first clear positive intervention on the 4-layer student, but the best student still remains far below the teacher. A small query-side projection head hurts retrieval quality, and the tested lightweight distillation objective does not materially improve the tradeoff. On the systems side, IVF is the only ANN variant that preserves enough quality to be worth discussion; HNSW is faster but too lossy under this setup.

These observations support a short, evidence-driven conclusion. Query-side compression is possible to study cleanly in an asymmetric setting, but shallow compression is highly brittle under simple truncation. Internal representation choices matter before broader efficiency claims can be made, and BEIR-style transfer checks \cite{thakur-etal-2021-beir} are useful only after the in-domain tradeoff is understood. The main value of the paper is therefore a controlled map of what fails first, what helps first, and which systems extension is plausible once the encoder story is stabilized.
